% Copyright (c) 2026 Ahmed Taha Bouallegui
% SPDX-License-Identifier: GPL-3.0-or-later

% =============================================================================
%  Ahmed Taha Bouallegui -- CV (version française)
%  Style : moderncv "classic" (the iconic name-banner + section-rule look)
%  Color : blue (matches personal branding)
%  Compile: pdflatex cv_bouallegui_fr.tex   (or xelatex)
%  Le contenu suit cv_bouallegui.tex : mêmes projets, mêmes dates, mêmes chiffres.
% =============================================================================
\documentclass[a4paper,11pt]{moderncv}
\moderncvstyle{classic}
\moderncvcolor{blue}

% Fonts / encoding
\usepackage[utf8]{inputenc}
\usepackage{lmodern}   % scalable Type1 fonts; microtype refuses to run with bitmap ones
\usepackage{graphicx}
\usepackage[french]{babel}

% Puces maitrisees : babel-french redefinit itemize (tiret demi-cadratin, grande marge
% gauche) apres tweaklist de moderncv. Un environnement "list" simple passe outre.
\newenvironment{cvlist}
  {\begin{list}{\strut\textcolor{color1}{\large\rmfamily\textbullet}}{\setlength{\leftmargin}{13pt}\setlength{\labelsep}{6pt}\setlength{\itemsep}{0pt}\setlength{\topsep}{1pt}\setlength{\parsep}{0pt}\setlength{\partopsep}{0pt}}}
  {\end{list}}

% Tighten geometry so everything fits on one page
\usepackage[top=7mm,bottom=5mm,left=12mm,right=12mm]{geometry}
\setlength{\parskip}{0pt}\setlength{\itemsep}{0pt}\setlength{\partopsep}{0pt}
% Un peu plus de place pour les dates, pour qu'elles tiennent sur une ligne.
\setlength{\hintscolumnwidth}{0.195\textwidth}

% Contact details
\name{Ahmed Taha}{Bouallegui}
\phone{+33 6 68 51 84 86}
\email{ahmedtbou@tutamail.com}
\address{Villeneuve d'Ascq, France}
\social[linkedin][www.linkedin.com/in/ahmed-taha-bouallegui-474215306]{LinkedIn}
\social[github]{TahaBouallegui}

% La photo d'identité est la convention française ; le CV anglais n'en a pas.
\photo[60pt][0.4pt]{images/cv/portrait-cv.jpg}

% Logo de l'école ou de l'entreprise, placé devant le nom de l'institution.
\newcommand{\cvlogo}[2][11pt]{\raisebox{-0.3em}{\includegraphics[height=#1]{#2}}\hspace{4pt}}

\begin{document}
\makecvtitle
\footnotesize

% ---------------------------------------------------------------------------
\section{Profil technique}
\cvitem[.05em]{}{Étudiant en ingénierie, spécialisé IoT et cybersécurité. Systèmes Linux, infrastructure déclarative avec Nix sur NixOS, dont un home lab qui fait tourner ses modèles en local. Je cherche un stage en administration systèmes ou en cybersécurité pour l'été 2027, 10 semaines à 4 mois.}

% ---------------------------------------------------------------------------
\section{Formation}
\cventry[.05em]{2023 à 2028}{M.Eng Internet of Things and Cybersecurity}{\cvlogo[13pt]{images/cv/polytech-mark.png}École Polytechnique Universitaire de Lille (Polytech Lille)}{Villeneuve d'Ascq, France}{}{}
\cventry[.05em]{2023 à 2025}{Cycle préparatoire (MPCSI)}{\cvlogo[13pt]{images/cv/polytech-mark.png}Polytech Lille}{Villeneuve d'Ascq, France}{}{Diplômé avec félicitations, major de promotion 2023/2024}
\cventry[.05em]{2023}{Baccalauréat}{France}{}{}{Mention Très Bien}

% ---------------------------------------------------------------------------
\section{Compétences techniques}
\cvitem[.05em]{Systèmes \& infra.}{Linux (NixOS, Debian, Arch, Ubuntu), KVM/QEMU, Docker, microVMs, systemd.}
\cvitem[.05em]{Config. déclarative}{Nix/NixOS, sops, builds reproductibles, gestion d'environnements.}
\cvitem[.05em]{IA locale}{llama.cpp, vLLM.}
\cvitem[.05em]{Réseau \& sécurité}{Tailscale, ZeroTier, isolation IAM, AWS (EC2, IAM, security groups), CTF (RE, web / sûreté mémoire).}
\cvitem[.05em]{Programmation}{Python, C/C++, C\#, Java, Rust, Bash, PostgreSQL, assembleur (AVR/ARM).}
\cvitem[.05em]{Outils}{Git, KiCad, MATLAB/Simulink, Matplotlib.}
\cvitem[.05em]{Langues}{anglais (C1), arabe (C2), français (C1), allemand (B1).}

% ---------------------------------------------------------------------------
\section{Projets sélectionnés}
\cventry[.05em]{Depuis août 2026}{Inférence LLM locale sur matériel grand public}{}{}{}{
\begin{cvlist}
  \item Modèles à poids ouverts sur trois machines avec \texttt{llama.cpp} : deux boîtiers RTX 3060 et un ThinkPad uniquement en CPU, atteignables seulement via un VPS.
  \item Un mélange d'experts de 125 Md de paramètres sur une carte de 12 Go : experts répartis entre VRAM, RAM et une table SSD de 29 Go. 35 tok/s en décodage et 253 en préfill, contre 12 et 55 avant le changement de moteur.
\end{cvlist}}

\cventry[.05em]{Depuis janv. 2026}{Infrastructure serveur Linux auto-hébergée}{}{}{}{
\begin{cvlist}
  \item Serveur headless : Immich, Jellyfin, Grocy, SearXNG et Samba. Commencé sous Debian en janvier 2026, reconstruit sous NixOS en février depuis un seul arbre de configuration.
  \item Rien n'écoute sur internet : les services sont atteignables uniquement via un VPS, et les secrets du dépôt sont chiffrés avec sops.
\end{cvlist}}

\cventry[.05em]{Janv. à juin 2026}{Clavier ergonomique split \textbf{(projet de cours)}}{}{}{}{Schéma et PCB deux couches dessinés dans KiCad, firmware écrit en C.}


% ---------------------------------------------------------------------------
\section{Expériences professionnelles}
\cventry[.05em]{Juin à juil. 2026}{Stage de recherche, CRIStAL \textbf{(équipe 2XS)}}{\cvlogo[13pt]{images/cv/univ-lille-mark.png}IRCICA \textbullet Université de Lille}{Lille, France}{}{
\begin{cvlist}
  \item Retiré chaque fork and exec de \texttt{iw} et \texttt{ip} du module wifi d'une sonde Linux de détection d'objets connectés au profit d'appels directs \texttt{nl80211} et rtnetlink, avec découverte automatique des interfaces.
  \item Contribué aux travaux de l'équipe 2XS sur l'empreinte mémoire, l'efficacité énergétique et la sécurité du logiciel embarqué.
\end{cvlist}}

\cventry[.05em]{Avr. 2025 à ce jour}{Fabmanager \textbf{(FabLab Fabricarium)}}{\cvlogo[13pt]{images/cv/polytech-mark.png}Polytech Lille}{Villeneuve d'Ascq, France}{}{
\begin{cvlist}
  \item Exploitation et maintenance du FabLab de l'école : imprimantes 3D, découpeuses laser, CNC et postes d'électronique.
  \item Accompagnement des étudiants qui prototypent du matériel et des projets IoT ; réservations, sécurité et sessions ouvertes.
\end{cvlist}}

\cventry[.05em]{Juil. à août 2023}{Stage en ingénierie de l'automatisation}{\cvlogo[10pt]{images/cv/valeo.png}Valeo}{Tunis, Tunisie}{}{
\begin{cvlist}
  \item Automatisation d'un contrôle qualité par reconnaissance sonore sur une ligne de production de volants, avec Python et Librosa.
\end{cvlist}}

\end{document}
